% Apéndice: construcción geométrica de las matrices de rotación 3D
\section{Construcción geométrica de las matrices de rotación 3D}\label{cap:rot3d}\label{ap:rot3d}

\subsection{Convención utilizada}

Se consideran vectores columna, rotaciones activas y un sistema cartesiano de mano derecha.
Una rotación activa mueve el vector mientras los ejes permanecen fijos. El sentido positivo
se determina mediante la regla de la mano derecha: el pulgar apunta en la dirección positiva
del eje y los dedos indican el sentido del giro.

\begin{idea}[Principio de construcción]
Una rotación alrededor de un eje:
\begin{itemize}
  \item deja fija la coordenada paralela al eje;
  \item aplica una rotación plana ordinaria en el plano perpendicular;
  \item tiene como columnas las imágenes de \(e_x,e_y,e_z\).
\end{itemize}
No es necesario memorizar tres matrices independientes: todas se reconstruyen a partir de
esta regla.
\end{idea}

\subsection{Recordatorio: rotación en un plano}

En un plano con coordenadas \((u,v)\), una rotación positiva de ángulo \(\theta\) cumple
\[
\begin{pmatrix}u'\\v'\end{pmatrix}
=
\begin{pmatrix}
\cos\theta&-\sin\theta\\
\sin\theta& \cos\theta
\end{pmatrix}
\begin{pmatrix}u\\v\end{pmatrix}.
\]
Por lo tanto,
\[
u'=u\cos\theta-v\sin\theta,
\qquad
v'=u\sin\theta+v\cos\theta.
\]
Para construir una rotación espacial basta identificar cuál es el par de coordenadas que
forma el plano perpendicular al eje.

\subsection{Rotación alrededor del eje \(z\)}

El eje \(z\) permanece fijo y la rotación ocurre en el plano \(xy\). Por lo tanto,
\[
\begin{aligned}
x'&=x\cos\theta-y\sin\theta,\\
y'&=x\sin\theta+y\cos\theta,\\
z'&=z.
\end{aligned}
\]
Agrupando los coeficientes de \(x,y,z\):
\[
\boxed{
R_z(\theta)=
\begin{pmatrix}
\cos\theta&-\sin\theta&0\\
\sin\theta& \cos\theta&0\\
0&0&1
\end{pmatrix}
}.
\]

La misma matriz se obtiene mirando las imágenes de la base:
\[
R_z e_x=
\begin{pmatrix}\cos\theta\\\sin\theta\\0\end{pmatrix},
\qquad
R_z e_y=
\begin{pmatrix}-\sin\theta\\\cos\theta\\0\end{pmatrix},
\qquad
R_z e_z=
\begin{pmatrix}0\\0\\1\end{pmatrix}.
\]
Estos tres vectores son, respectivamente, las tres columnas de \(R_z\).

\begin{center}
\begin{tikzpicture}[scale=1.0]
  \draw[axisline,-{Latex[length=2.4mm]}] (-.3,0)--(5.3,0) node[right] {$x$};
  \draw[axisline,-{Latex[length=2.4mm]}] (0,-.3)--(0,4.5) node[above] {$y$};
  \draw[mainline,-{Latex[length=2.7mm]}] (0,0)--(4,0)
    node[below right] {$e_x$};
  \draw[rust,line width=1.05pt,-{Latex[length=2.7mm]}] (0,0)--(3,2.6)
    node[above right] {$R_ze_x$};
  \draw[teal,line width=.8pt,-{Latex[length=2mm]}] (1.4,0)
    arc[start angle=0,end angle=41,radius=1.4];
  \node[text=teal] at (1.7,.55) {$\theta$};
  \node[font=\small,text=muted,align=center] at (7,2.0)
    {\(z\) sale perpendicular al dibujo\\y permanece fijo};
\end{tikzpicture}
\end{center}

\subsection{Rotación alrededor del eje \(x\)}

Ahora \(x\) permanece constante y la rotación ocurre en el plano \(yz\). Se aplica la misma
matriz plana al par \((y,z)\):
\[
\begin{aligned}
x'&=x,\\
y'&=y\cos\theta-z\sin\theta,\\
z'&=y\sin\theta+z\cos\theta.
\end{aligned}
\]
Por lo tanto,
\[
\boxed{
R_x(\theta)=
\begin{pmatrix}
1&0&0\\
0&\cos\theta&-\sin\theta\\
0&\sin\theta& \cos\theta
\end{pmatrix}
}.
\]
Sus columnas vuelven a ser las imágenes de la base:
\[
R_xe_x=e_x,\qquad
R_xe_y=
\begin{pmatrix}0\\\cos\theta\\\sin\theta\end{pmatrix},
\qquad
R_xe_z=
\begin{pmatrix}0\\-\sin\theta\\\cos\theta\end{pmatrix}.
\]
La primera columna \((1,0,0)^T\) expresa geométricamente que el eje \(x\) permanece fijo.

\subsection{Rotación alrededor del eje \(y\)}

En este caso \(y\) permanece constante y el giro ocurre en el plano \(xz\). La regla de la
mano derecha determina los signos:
\[
\begin{aligned}
x'&=x\cos\theta+z\sin\theta,\\
y'&=y,\\
z'&=-x\sin\theta+z\cos\theta.
\end{aligned}
\]
Así:
\[
\boxed{
R_y(\theta)=
\begin{pmatrix}
\cos\theta&0&\sin\theta\\
0&1&0\\
-\sin\theta&0&\cos\theta
\end{pmatrix}
}.
\]
Por columnas:
\[
R_ye_x=
\begin{pmatrix}\cos\theta\\0\\-\sin\theta\end{pmatrix},
\qquad
R_ye_y=e_y,
\qquad
R_ye_z=
\begin{pmatrix}\sin\theta\\0\\\cos\theta\end{pmatrix}.
\]

\begin{atencion}[Por qué \(R_y\) parece tener signos diferentes]
Si se escribe directamente la rotación 2D sobre el par ordenado \((x,z)\), puede aparecer
una confusión de orientación. Mirado desde el extremo positivo del eje \(y\), el plano
\(xz\) debe orientarse de acuerdo con la regla de la mano derecha. Una comprobación segura es
usar un ángulo de \(90^\circ\): la matriz anterior lleva \(e_x\) a \(-e_z\) y \(e_z\) a
\(e_x\).
\end{atencion}

\subsection{Resumen geométrico}

\begin{center}
\begin{tabularx}{.9\textwidth}{>{\bfseries}c c c}
\toprule
Rotación & Coordenada fija & Plano que rota\\
\midrule
\(R_x\) & \(x\) & \(yz\)\\
\(R_y\) & \(y\) & \(zx\), orientado por mano derecha\\
\(R_z\) & \(z\) & \(xy\)\\
\bottomrule
\end{tabularx}
\end{center}

Cada matriz debe satisfacer cuatro comprobaciones inmediatas:
\[
R_i(0)=I,\qquad
R_i(\theta)^TR_i(\theta)=I,\qquad
\det R_i(\theta)=1,\qquad
R_i(\theta)e_i=e_i.
\]
La última igualdad expresa que el vector director del eje es un autovector con autovalor
\(1\).

\subsection{Rotación alrededor de un eje arbitrario}

Sea \(u\) un vector unitario que define el eje. Todo vector \(v\) puede separarse en una parte
paralela y otra perpendicular:
\[
v=v_\parallel+v_\perp,
\qquad
v_\parallel=(u\cdot v)u,
\qquad
v_\perp=v-(u\cdot v)u.
\]
La rotación deja fija \(v_\parallel\). En el plano perpendicular, \(v_\perp\) gira como en
dos dimensiones y \(u\times v\) proporciona la dirección perpendicular dentro de ese plano:
\[
R(u,\theta)v
=
v_\parallel
+\cos\theta\,v_\perp
+\sin\theta\,(u\times v).
\]
Reagrupando:
\[
\boxed{
R(u,\theta)
=
\cos\theta\,I
+(1-\cos\theta)uu^T
+\sin\theta\,[u]_\times
},
\]
donde
\[
[u]_\times=
\begin{pmatrix}
0&-u_z&u_y\\
u_z&0&-u_x\\
-u_y&u_x&0
\end{pmatrix}.
\]
Esta es la fórmula de Rodrigues. Las matrices \(R_x,R_y,R_z\) se recuperan tomando
\(u=e_x,e_y,e_z\), respectivamente.

\begin{resultado}[Idea central]
Una matriz de rotación 3D se construye dejando fija la componente paralela al eje y rotando
la componente perpendicular como en un plano. Sus columnas muestran dónde terminan los tres
vectores de la base.
\end{resultado}

